\documentclass[12pt]{article}
\usepackage[T1]{fontenc}
\usepackage[utf8]{inputenc}
\usepackage{lmodern}
\usepackage{textcomp}
\usepackage{enumitem}
\usepackage{geometry}
\geometry{margin=1in}
\begin{document}
\begin{center}
Answer Key - Economics - K-12 Level Assessment 19 \
\small (Answers are ordered exactly as in the question paper.)
\end{center}

\section*{Multiple Choice}
\begin{enumerate}[label=\arabic*.]
\item \textbf{Answer:} B
\\ \textit{Options:} A) inflation had already impacted the economy before the fiscal stimulus.; B) the economy initially had some unemployed resources.; C) aggregate supply decreased.; D) aggregate demand is steeply sloped.
\\ \textit{Note:} Correct option: B - the economy initially had some unemployed resources.
\item \textbf{Answer:} C
\\ \textit{Options:} A) increase interest rates and throw the economy into a recession.; B) increase interest rates and depreciate the nation's currency.; C) decrease interest rates and risk an inflationary period.; D) decrease interest rates and throw the economy into a recession.
\\ \textit{Note:} Correct option: C - decrease interest rates and risk an inflationary period.
\item \textbf{Answer:} D
\\ \textit{Options:} A) the price of a good divided by its marginal utility.; B) the marginal utility of the good divided by its price.; C) the total utility of the good.; D) the difference between the consumer's value and the market price.
\\ \textit{Note:} Correct option: D - the difference between the consumer's value and the market price.
\item \textbf{Answer:} A
\\ \textit{Options:} A) a monopsony buys from a monopoly; B) a monopoly sells to two different types of consumers; C) a monopoly buys from a monopsony; D) a monopolist sells two different types of goods
\\ \textit{Note:} Correct option: A - a monopsony buys from a monopoly
\item \textbf{Answer:} B
\\ \textit{Options:} A) nominal GDP falls.; B) the nominal interest rate falls.; C) bond prices fall.; D) the supply of money falls.
\\ \textit{Note:} Correct option: B - the nominal interest rate falls.
\end{enumerate}

\section*{True / False}
\begin{enumerate}[label=\arabic*.]
\setcounter{enumi}{5}
\item \textbf{Answer:} True
\\ \textit{Options:} A) True; B) False
\\ \textit{Note:} LLM-generated true statement based on MCQ.
\item \textbf{Answer:} True
\\ \textit{Options:} A) True; B) False
\\ \textit{Note:} LLM-generated true statement based on MCQ.
\item \textbf{Answer:} True
\\ \textit{Options:} A) True; B) False
\\ \textit{Note:} LLM-generated true statement based on MCQ.
\item \textbf{Answer:} False
\\ \textit{Options:} A) True; B) False
\\ \textit{Note:} LLM-generated false statement. Correct answer: 'Savers'.
\item \textbf{Answer:} True
\\ \textit{Options:} A) True; B) False
\\ \textit{Note:} LLM-generated true statement based on MCQ.
\end{enumerate}

\section*{Long Form Response}
\begin{enumerate}[label=\arabic*.]
\setcounter{enumi}{10}
\item \textbf{Answer:} Increases Increases Decreases. Explain why this is the correct answer and provide supporting evidence. Reference key facts or concepts that support this choice.
\\ \textit{Note:} Long-form derived from MCQ; award credit for stating the correct idea and giving supporting reasoning.
\item \textbf{Answer:} Substitution effects and income effects. Explain why this is the correct answer and provide supporting evidence. Reference key facts or concepts that support this choice.
\\ \textit{Note:} Long-form derived from MCQ; award credit for stating the correct idea and giving supporting reasoning.
\end{enumerate}

\vfill
\noindent\textit{End of Answer Key}
\end{document}